\section*{第 \ref{ch:b1:flowering-plant-organization} 章——开花植物的功能组织}

\begin{solution}{exo:b1:flowering-plant-organization:1}
根：固定植株，吸收水分和无机离子。茎：支持
叶片向光伸展，并在根与叶之间输导物质。
叶：光合作用（以及气体交换、蒸腾作用）。
\end{solution}

\begin{solution}{exo:b1:flowering-plant-organization:2}
模块指一个节点连同其上的叶和腋芽，以及下方
的\omterm{def:b1:flowering-plant-organization:organs}{节间}。\omterm{prop:b1:flowering-plant-organization:modular}{无限生长的}：没有固定的成体大小；顶芽和腋芽
只要植株存活，就不断增加新的模块。
\end{solution}

\begin{solution}{exo:b1:flowering-plant-organization:3}
\omterm{def:b1:flowering-plant-organization:tissues}{表皮}：活的；起保护作用，具\omterm{def:b1:flowering-plant-organization:tissues}{角质层}、\omterm{def:b1:flowering-plant-organization:tissues}{气孔}和\omterm{def:b1:flowering-plant-organization:organs}{根毛}。\omterm{def:b1:flowering-plant-organization:tissues}{薄壁组织}：
活的；进行光合作用，贮藏物质。\omterm{def:b1:flowering-plant-organization:tissues}{厚角组织}：活的；提供柔韧支持。
\omterm{def:b1:flowering-plant-organization:tissues}{厚壁组织}：死的；提供刚性支持。\omterm{def:b1:flowering-plant-organization:tissues}{木质部}\omterm{def:b1:flowering-plant-organization:tissues}{导管}：死的；
水分运输。\omterm{def:b1:flowering-plant-organization:tissues}{韧皮部}\omterm{def:b1:flowering-plant-organization:tissues}{筛管}：活的（有伴胞）；糖类
运输。
\end{solution}

\begin{solution}{exo:b1:flowering-plant-organization:4}
$e^{-0.9} = 0.41$，即\ \qty{41}{\%}；$e^{-3} = 0.05$，即\
\qty{5}{\%}。
\end{solution}

\begin{solution}{exo:b1:flowering-plant-organization:5}
双子叶植物的根（单一中柱、\omterm{prop:b1:flowering-plant-organization:stemroot}{内皮层}，\omterm{def:b1:flowering-plant-organization:tissues}{木质部}与\omterm{def:b1:flowering-plant-organization:tissues}{韧皮部}
沿半径交替排列；\omterm{def:b1:flowering-plant-organization:tissues}{木质部}脊数少，无中央髓）。
\end{solution}

\begin{solution}{exo:b1:flowering-plant-organization:6}
弯曲的梁，应力在表面最大，沿轴线完全不受应力，
故置于外周环的材料抵抗弯曲时质量最省
（即管状）。受拉缆绳的截面应力均匀，
可自由弯曲；中央单股承受拉力，而皮层保持柔韧，
侧根可从中穿出。
\end{solution}

\begin{solution}{exo:b1:flowering-plant-organization:7}
$1 - e^{-2.4} = 0.91$；当 $L = 6$ 时，$1 - e^{-3.6} = 0.97$：
再加两层叶，可多截获 \qty{6}{\%} 的光，
却需多投入 \qty{50}{\%} 的叶量。
\end{solution}

\begin{solution}{exo:b1:flowering-plant-organization:8}
$S = \pi d L$：$d = 240/(\pi\times 623\,000) = \qty{1.2e-4}{m}$，约
\qty{0.12}{mm}。$(240 + 400)/5 = 128$。
\end{solution}

\begin{solution}{exo:b1:flowering-plant-organization:9}
每毫米根的表面积：$\pi\times 0.4 = \qty{1.26}{mm^2}$。每毫米\omterm{def:b1:flowering-plant-organization:organs}{根毛}：
$250\times\pi\times 0.010\times 0.8 = \qty{6.3}{mm^2}$。总计
$7.5/1.26 = 6$：\omterm{def:b1:flowering-plant-organization:organs}{根毛}使表面积增大为六倍。
\end{solution}

\begin{solution}{exo:b1:flowering-plant-organization:10}
在林下，较低处的腋芽受遮荫，它们
所能形成的模块难以自给，因此它们保持休眠，或其
分枝死亡，生长则分配给向上伸向光的顶芽，
在开阔处，每个芽都受光，每个分枝都有回报，
因此树冠在低处也能充实。代价是：细高的树干
是对无产出木材的大量投入，且易受风害；收益是：植株
使其形态与光照分布相适应，而无需任何固定的规划。
\end{solution}

\begin{solution}{exo:b1:flowering-plant-organization:11}
\omterm{def:b1:flowering-plant-organization:tissues}{气孔}窝和绒毛在\omterm{def:b1:flowering-plant-organization:tissues}{气孔}上方保持一层静止、潮湿的空气；
从气室到流动空气的蒸气梯度分布在更长的路径上，
蒸腾作用随之下降。二氧化碳必须沿同样更长的路径扩散进入，
因此吸收也下降，但降幅较小，因为
$\mathrm{CO_2}$ 梯度由大气和叶的消耗共同决定，
在沙漠中水是限制性资源，这一权衡
挽救了植物的生命；在雨林中水分充足，碳的获取
才是竞争所奖励的。
\end{solution}

\begin{solution}{exo:b1:flowering-plant-organization:12}
植物没有受调节的\omterm{def:b1:mammal-organization:compartments}{细胞外液}，使细胞与外界之间
的温度和成分保持恒定：其
细胞直接接触土壤水、空气和阳光，温度与空气相同。
因此每个细胞自行调节其内含物（离子、
水分、pH）穿过其膜与细胞壁；植物则在表面调节
交换（\omterm{def:b1:flowering-plant-organization:tissues}{气孔}开闭，\omterm{prop:b1:flowering-plant-organization:stemroot}{内皮层}选择进入\omterm{def:b1:flowering-plant-organization:tissues}{木质部}的物质），
并通过生长调整体态（干燥土壤中多长根，
遮荫处多长叶）。它确实进行调节，但
调节发生在表面与形态上，而非专属的体液中。
\end{solution}

\begin{solution}{pb:b1:flowering-plant-organization:1}
\textbf{1.} $5\times 100 = \qty{500}{m^2}$。
\textbf{2.} $e^{-2.5} = 0.082$ 到达地面，\qty{92}{\%}
被截获。
\textbf{3.} $L = 2$：$1 - e^{-1} = \qty{63}{\%}$；$L = 8$：$1 -
e^{-4} = \qty{98}{\%}$。
\textbf{4.} 从 2 增至 5 时，树冠以 \qty{300}{m^2} 叶面积
换得 \qty{29}{\%} 的光；从 5 增至 8 时，
另外 \qty{300}{m^2} 叶面积只能多得 \qty{6}{\%} 的光，
而其建造与供水成本相当：多出的叶子不划算。
\textbf{5.} $500/0.005 = \num{100000}$ 片叶。
\textbf{6.} $500\times 10^6\,\mathrm{mm^2}\times 150 = \num{7.5e10}$
\omterm{def:b1:flowering-plant-organization:tissues}{气孔}。
\textbf{7.} $\qty{50}{\micro m^2}$；每 mm$^2$：$150\times 50\times
10^{-6} = \qty{7.5e-3}{mm^2}$，即\ \qty{0.75}{\%} 的表面积。
\textbf{8.} $500\times 4\times 10^{-6} = \qty{2e-3}{mol/s}$，即\
\qty{88}{mg/s}。
\textbf{9.} $0.088\times 43\,200 = \qty{3.8}{kg}$ 的 $\mathrm{CO_2}$，
其中含碳 $3.8\times 12/44 = \qty{1.04}{kg}$。
\textbf{10.} $150\times 1.04 = \qty{156}{kg}$ 的碳；一半，
\qty{78}{kg} 进入木材，按含碳 \qty{50}{\%} 计，折合 \qty{156}{kg} 木材。
\textbf{11.} $500\times 10^{-3} = \qty{0.5}{mol/s}$，即\ \qty{9}{g/s}。
\textbf{12.} $9\times 43\,200 = \qty{389}{kg}$，约 \qty{390}{L}。
\textbf{13.} $0.5/(2\times 10^{-3}) = 250$ 个水分子，每
$\mathrm{CO_2}$。让 $\mathrm{CO_2}$ 进入的\omterm{def:b1:flowering-plant-organization:tissues}{气孔}也让水分逸出；
\omterm{def:b1:flowering-plant-organization:tissues}{气孔}内空气已饱和，而 $\mathrm{CO_2}$ 仅占
外界空气的 \qty{0.04}{\%}，因此向外的蒸气梯度
远比向内的 $\mathrm{CO_2}$ 梯度陡峭。
\textbf{14.} $\qty{2}{mm}\times\qty{100}{m^2} = \qty{200}{L}$：相当于
半天的量。
\textbf{15.} \qty{100}{m^3}；$5\times 100 = \qty{500}{km}$。
\textbf{16.} $\pi\times 0.5\times 10^{-3}\times 5\times 10^5 =
\qty{785}{m^2}$。
\textbf{17.} $200\times 5\times 10^8\,\mathrm{mm} = \num{1e11}$ \omterm{def:b1:flowering-plant-organization:organs}{根毛}。
\textbf{18.} 每根 $\pi\times 10^{-5}\times 5\times 10^{-4} =
\qty{1.57e-8}{m^2}$；合计 \qty{1570}{m^2}。
\textbf{19.} $785 + 1570 = \qty{2355}{m^2}$，为叶面积的 $4.7$ 倍。
\textbf{20.} \qty{1}{mm} 以内的体积：$\pi\times(10^{-3})^2\times 5
\times 10^5 = \qty{1.6}{m^3}$，即占土壤的\ \qty{1.6}{\%}。
\textbf{21.} 黑麦：$\pi\times 10^{-6}\times 6.2\times 10^5 =
\qty{1.95}{m^3}$，是 \qty{0.05}{m^3} 箱体的 \qty{3900}{\%}：每
土壤中每一点都在一毫米范围内有四十条根。黑麦
探索土壤的彻底程度高出两千多倍 ——
一种仅有四个月生长期的作物，对比一棵探索一百
立方米土壤的树。
\textbf{22.} $1000 + 2355 = \qty{3355}{m^2}$；每平方米地面
$\qty{34}{m^2}$。
\textbf{23.} 人：$130/2 = 65$；橡树的比值为人类
折叠表面的一半，且无需任何折叠。
\textbf{24.} 橡树每年增加模块、叶和根，因此其
表面随年龄增长；哺乳动物的表面是固定的\omterm{def:b1:mammal-organization:organ}{器官}，
成年时即已定型，提高比值的唯一办法是在发育中
将其表面折叠起来。
\textbf{25.} 每平方米地面约有 \qty{34}{m^2} 的交换表面：
地上为 \qty{10}{m^2} 叶片（两面），地下为 \qty{24}{m^2} 的
根和\omterm{def:b1:flowering-plant-organization:organs}{根毛}。
\end{solution}
